% Gliederung Stand 30.10.2026. Grundlagen + Stand der Technik, ~8 Seiten.
\chapter{Simulation und Software-Infrastruktur}
\label{chap:simulation}
% Kurze Kapiteleinleitung vor 3.1.

\section{Simulation in der Robotik}
\label{sec:sim_robotik}

\subsection{Einsatzzwecke der Simulation}
\label{sec:sim_zwecke}
% Entwicklung, virtuelle Inbetriebnahme, Schulung. Kandidat: Choi2021.

\subsection{Sim-to-Real-Lücke}
\label{sec:sim2real}
% Ursachen und Ansätze zur Verringerung. Kandidaten: Zhao2020, Tobin2017, Tan2018.

\section{Physiksimulatoren}
\label{sec:physiksimulatoren}

\subsection{MuJoCo}
\label{sec:mujoco}
% Kontaktmodell, Integratoren, Zeitschritt. Kandidaten: Todorov2012, Todorov2014, Tunyasuvunakool2020.

\subsection{Vergleich mit weiteren Simulatoren}
\label{sec:sim_vergleich}
% Isaac Sim, Gazebo, PyBullet. Kandidaten: Erez2015, Collins2021.

\section{Middleware und Bereitstellung}
\label{sec:middleware}
% ROS 2, DDS-Domains, Container mit Docker.

\section{Sicherheit beim Betrieb von Robotern}
\label{sec:sicherheit_grundlagen}
% Normen (z. B. ISO 10218-1), Not-Halt, Zustimmeinrichtung (Totmann), Geschwindigkeitsbegrenzung.
